% ============================================================================
% Fichier  : partie1/P01_C04_fondements_theoriques.tex
% Titre    : Fondements Theoriques
% Auteur   : MINKA MI NGUIDJOÏ Thierry Emmanuel
% Version  : 1.0 -- Juin 2026
% Statut   : REDIGÉ
% Sources  : 3_fondements_theoriques.tex, 21_fondements_cryptanalyse.tex,
%             4_etat_art.tex, minka2025cro, minka2025zknr
% ============================================================================

\chapter{Fondements Théoriques}
\label{chap:fond-theo}

\epigraph{%
    « La théorie sans la pratique est vaine~; la pratique sans la théorie
    est aveugle. L'investigation numérique exige les deux. »%
}{--- Adaptation d'Emmanuel Kant, \textit{Critique de la raison pure}}

\niveauChapitre{Fondamental}{%
    Mathématiques de niveau licence (probabilités, algèbre de base,
    notions de complexité). Les résultats sont commentés pour être
    accessibles sans démonstration exhaustive.%
}

\begin{boxobjectifs}
\begin{itemize}
    \item Maîtriser le \termecle{principe de Locard numérique} et
          sa taxonomie des traces~: primaires, secondaires, tertiaires,
          résiduelles.
    \item Appliquer les trois outils de la théorie de l'information~:
          entropie, distance de Hamming, distance de Levenshtein.
    \item Modéliser un réseau criminel avec la théorie des graphes
          (centralité de degré, d'intermédiarité, PageRank, graphes temporels).
    \item Comprendre le modèle Dolev-Yao et en déduire les menaces
          sur la preuve numérique.
    \item Situer les modèles de processus (DFRWS, Casey, NIST, ISO~27037)
          dans leur cohérence et leurs différences.
    \item Articuler fondements théoriques et Trilemme CRO.
\end{itemize}
\end{boxobjectifs}

\bigskip

% ============================================================================
\section{Le Principe de Locard Numérique}
\label{sec:locard}
% ============================================================================

\subsection{Origine et transposition au numérique}

Édmond Locard (1877--1966), directeur du premier laboratoire de police
scientifique au monde (Lyon, 1910), a formulé le principe fondateur de la
criminalistique~: \textit{« Tout contact laisse une trace. »} Sa transposition
au monde numérique est non triviale mais inévitable.

\begin{boxdef}
\textbf{Principe de Locard numérique}\index{Locard, principe numérique} ---
Tout processus informatique laisse des traces dans son environnement
d'exécution. Ces traces sont automatiques, involontaires, et persistent au
sein des systèmes de fichiers, des journaux, de la mémoire et du réseau,
indépendamment de la volonté de l'acteur qui a engendré le processus.
\end{boxdef}

La différence cruciale avec le monde physique est que le numérique laisse
des traces \textbf{à plusieurs niveaux simultanément}~: un simple clic sur
un lien Web génère des entrées dans l'historique du navigateur, une requête
DNS dans les journaux réseau, une connexion TCP dans les logs du routeur,
un artefact de cache dans le répertoire temporaire, et potentiellement
un enregistrement dans les bases de données du serveur distant.

\subsection{Taxonomie des traces numériques}

\begin{tableaucours}{p{3.2cm}p{4.2cm}p{4.8cm}}{%
    \tableauHeader{Catégorie} &
    \tableauHeader{Définition} &
    \tableauHeader{Exemples concrets}
}
\textbf{Traces primaires} &
    Générées directement par l'action incriminée &
    Fichier créé, email envoyé, transaction MoMo exécutée, connexion RDP \\
\textbf{Traces secondaires} &
    Générées par l'environnement en réponse à l'action &
    Logs système, historique navigateur, registre Windows, Prefetch,
    métadonnées EXIF \\
\textbf{Traces tertiaires} &
    Générées par les outils de sécurité &
    Alertes IDS/SIEM, logs pare-feu, journaux VPN, audit Active Directory \\
\textbf{Traces résiduelles} &
    Persistant après tentative d'effacement &
    Fichiers dans l'espace non alloué, entrées MFT orphelines,
    \textit{slack space}, \texttt{\$LogFile} NTFS \\
\end{tableaucours}
\vspace{-0.8em}
\begin{center}{\small\itshape Tableau~4.1 --- Taxonomie des traces numériques
selon le principe de Locard}\end{center}

\subsection{Locard dans les affaires étudiées}

Le principe de Locard explique a posteriori pourquoi chaque grande affaire
a pu être résolue~:
\begin{itemize}
    \item \textbf{BTK/Rader}~: les métadonnées Office sont des \emph{traces
          secondaires} --- générées automatiquement par Word, indépendamment
          de toute intention de Rader.
    \item \textbf{MVOM}~: la connexion RDP est une \emph{trace primaire} dans
          les logs pare-feu~; le processus \texttt{svchost.exe} malveillant en
          RAM est une \emph{trace secondaire} de son installation.
    \item \textbf{WOURI}~: les transactions Bitcoin sont des \emph{traces
          primaires} permanentes et immuables --- Locard dans sa forme absolue.
\end{itemize}

\begin{boiteRemarque}{L'anti-forensique comme tentative de violation de Locard}
Les techniques d'\termecle{anti-forensique}\index{anti-forensique} (suppression
de fichiers, nettoyage des logs, chiffrement) ne violent pas Locard~: elles
tentent de le contourner. Mais l'effacement lui-même laisse des traces~:
les outils utilisés apparaissent dans le Prefetch Windows~; les entrées MFT
supprimées persistent dans l'espace non alloué~; l'horodatage de la suppression
est enregistré dans \texttt{\$LogFile}. La trace de la suppression est une
trace secondaire au sens de Locard.
\end{boiteRemarque}

% ============================================================================
\section{Théorie de l'Information~: l'Arsenal de Shannon}
\label{sec:shannon-arsenal}
% ============================================================================

\subsection{L'entropie comme détecteur d'anomalie}

La théorie de Shannon \cite{shannon1948} fournit trois outils directement
opérationnels. L'entropie $H(X)$ a été introduite au Chapitre~\ref{chap:phi-fond}.
Rappelons son calcul sur un fichier binaire~: on estime la distribution des
256 valeurs d'octet possibles~:

\begin{equation}
H_{\text{fichier}} = -\sum_{b=0}^{255}
\frac{n_b}{N}\,\log_2\!\left(\frac{n_b}{N}\right)
\label{eq:shannon-fichier}
\end{equation}

où $n_b$ est le nombre d'occurrences de l'octet de valeur $b$ et $N$ le nombre
total d'octets analysés.

\begin{lstlisting}[language=bash_forensique,
    caption={Calcul de l'entropie d'un fichier --- script Python forensique}]
import math, collections

def entropie(chemin, echantillon=65536):
    with open(chemin, 'rb') as f:
        d = f.read(echantillon)
    n = len(d)
    cpt = collections.Counter(d)
    h = -sum((c/n)*math.log2(c/n) for c in cpt.values())
    return h

# Interprétation
h = entropie('suspect.bin')
if   h > 7.5: print('CHIFFRE ou COMPRESSE -- investigation approfondie')
elif h > 6.0: print('Binaire compile -- normal pour un executable')
elif h > 3.5: print('Texte ou donnees structurees -- normal')
else        : print('Tres repetitif -- fichier sparse ou quasi vide')
\end{lstlisting}

\subsection{La distance de Hamming~: identifier les variantes de malware}

La \termecle{distance de Hamming}\index{distance de Hamming} mesure le nombre
de positions où deux chaînes binaires de même longueur diffèrent~:

\begin{equation}
d_H(x, y) = \sum_{i=1}^{n} \mathbf{1}[x_i \neq y_i]
\label{eq:hamming}
\end{equation}

En pratique, les \termecle{fuzzy hashes}\index{fuzzy hash} (ssdeep, TLSH)
calculent des distances approximatives sur des fichiers de taille variable~:
une distance normalisée $d_H(x,y)/n < 0.15$ entre deux binaires indique
une probable appartenance à la même famille de malware.

\begin{lstlisting}[language=bash_forensique,
    caption={Fuzzy hashing avec ssdeep}]
# Calculer le fuzzy hash d'un fichier
ssdeep -b malware_suspect.exe

# Comparer avec une référence (score 0-100, 100 = identiques)
ssdeep -m malware_reference.exe malware_suspect.exe
# Sortie : matches malware_reference.exe (88)
# Score 88 : variante très probable de la même famille

# Comparaison sur un répertoire entier
ssdeep -r -m malware_reference.exe /dossier/suspects/
\end{lstlisting}

\subsection{La distance de Levenshtein~: typosquatting et comparaison de textes}

La \termecle{distance de Levenshtein}\index{distance de Levenshtein} mesure
le nombre minimum d'insertions, suppressions et substitutions pour transformer
une chaîne en une autre~:

\begin{equation}
d_L(s, t) = \min_{\text{séq. d'opérations}} |\text{opérations}(s \to t)|
\label{eq:levenshtein}
\end{equation}

Applications forensiques directes~:
\begin{itemize}
    \item \textbf{Typosquatting} (WOURI)~: le domaine frauduleux
          \texttt{medequip-portugal.com} a une distance de Levenshtein très faible
          par rapport au domaine légitime \texttt{medequip.pt}~; l'outil
          \textit{dnstwist} automatise cette détection.
    \item \textbf{Variantes de scripts}~: deux scripts avec $d_L$ faible ont
          probablement le même auteur.
    \item \textbf{Pseudonymes}~: deux noms très proches dans des bases de données
          différentes peuvent désigner la même personne.
\end{itemize}

\begin{lstlisting}[language=bash_forensique,
    caption={Détection de typosquatting avec dnstwist}]
# Génère toutes les variantes du domaine et vérifie celles enregistrées
dnstwist --registered medequip.pt

# Sortie typique :
# Original*       medequip.pt           185.xx.xx.xx
# Hyphenation     medequip-pt.com       91.xxx.xx.xx  <- SUSPECT
# Transposition   meidquip.pt           [NON ENREGISTRE]
\end{lstlisting}

% ============================================================================
\section{Théorie des Graphes en Investigation Numérique}
\label{sec:graphes}
% ============================================================================

\subsection{Formalisation des réseaux criminels}

Un réseau criminel ou de communication est modélisé par un
\termecle{graphe orienté pondéré}\index{graphe!orienté pondéré}~:

\begin{equation}
G = (V, E, w), \quad
E \subseteq V \times V, \quad
w : E \to \mathbb{R}^+
\label{eq:graphe-pondere}
\end{equation}

où $V$ est l'ensemble des entités (personnes, machines, comptes, adresses IP),
$E$ l'ensemble des relations (communications, transferts, accès) et $w(u,v)$
le poids de la relation (montant transféré, fréquence d'accès).

\subsection{Métriques de centralité~: identifier les acteurs clés}

\subsubsection{Centralité de degré}

\begin{equation}
C_D(v) = \frac{\deg(v)}{|V| - 1}
\label{eq:centralite-degre}
\end{equation}

Un nœud à forte centralité de degré est un \textbf{hub de communication}.
Dans WOURI, ATANGANA avait la centralité de degré la plus élevée.

\subsubsection{Centralité d'intermédiarité (\textit{betweenness centrality})}

\begin{equation}
C_B(v) = \sum_{s \neq v \neq t}
\frac{\sigma(s, t \mid v)}{\sigma(s, t)}
\label{eq:centralite-betweenness}
\end{equation}

où $\sigma(s,t)$ est le nombre de chemins les plus courts entre $s$ et $t$,
et $\sigma(s,t \mid v)$ ceux qui passent par $v$. Un acteur à forte centralité
d'intermédiarité est un \textbf{pont critique}~: le neutraliser fragmente
le réseau. Dans WOURI, NGUEMO était le pont entre la collecte des fonds
(Cameroun) et le blanchiment (EAU).

\subsubsection{PageRank forensique}

\begin{equation}
\mathrm{PR}(v) = \frac{1-d}{|V|} + d \sum_{u \in \mathcal{N}^{-}(v)}
\frac{\mathrm{PR}(u)}{\mathrm{deg}^+(u)}, \quad d \approx 0.85
\label{eq:pagerank}
\end{equation}

En investigation d'emails, PageRank identifie les personnes les plus influentes
dans un réseau de communication~: celles dont les messages sont fréquemment
répondus ou transférés (comme dans l'affaire Enron).

\subsection{Graphes temporels~: la chronologie comme dimension analytique}

Un \termecle{graphe temporel}\index{graphe!temporel} $G = (V, E, T)$ ajoute
une dimension temporelle aux arêtes~:

\begin{equation}
E \subseteq V \times V \times T
\label{eq:graphe-temporel}
\end{equation}

En forensique réseau, chaque connexion est un triplet
$(\text{source}, \text{destination}, t)$. L'analyse du graphe temporel
permet de \textbf{reconstituer un chemin d'attaque}~: deux événements $e_1$
et $e_2$ sont causalement liés si $t(e_2) - t(e_1)$ est faible et si les
entités impliquées se recoupent.

\begin{lstlisting}[language=bash_forensique,
    caption={Analyse de graphe criminel avec NetworkX}]
import networkx as nx

G = nx.DiGraph()
# Arêtes : (source, destination, poids)
edges = [
    ('ATANGANA', 'OKAFOR',  {'w': 1, 'd': '2026-01'}),
    ('OKAFOR',   'SOTRADI', {'w': 47, 'd': '2026-01-15'}),
    ('SOTRADI',  'NGUEMO',  {'w': 47, 'd': '2026-01-15'}),
    ('NGUEMO',   'HASSAN',  {'w': 45, 'd': '2026-01-15'}),
]
G.add_edges_from([(u,v,d) for u,v,d in edges])

# Centralité d'intermédiarité
btw = nx.betweenness_centrality(G)
# -> NGUEMO et OKAFOR ont les valeurs les plus élevées

# PageRank (acteur le plus 'central' dans le flux)
pr = nx.pagerank(G, weight='w')
\end{lstlisting}

% ============================================================================
\section{Le Modèle Dolev-Yao~: l'Adversaire Formel}
\label{sec:dolev-yao}
% ============================================================================

\subsection{Définition}

Le \termecle{modèle Dolev-Yao}\index{Dolev-Yao, modèle} \cite{dolevyao1983}
est le cadre formel standard pour raisonner sur la sécurité d'un protocole
cryptographique face à un adversaire actif~:

\begin{boxdef}
\textbf{Adversaire Dolev-Yao} --- Un adversaire $\mathcal{A}$ qui~:
\begin{enumerate}
    \item peut \textbf{intercepter} tout message transitant sur le réseau~;
    \item peut \textbf{modifier, rejouer, supprimer ou réordonner} tout message~;
    \item peut \textbf{injecter} des messages en se faisant passer pour
          n'importe quelle entité~;
    \item \textbf{ne peut pas} casser les primitives cryptographiques
          correctement utilisées (déchiffrer sans clé, forger une signature).
\end{enumerate}
\end{boxdef}

\subsection{Implications forensiques}

\textbf{Implication~1~:} toute preuve transmise sur un réseau non sécurisé
est potentiellement compromise. La chaîne de custody numérique doit inclure
la sécurisation des canaux de transmission (chiffrement + signature).

\textbf{Implication~2~:} la non-répudiation est un problème difficile.
Comment prouver qu'un message a bien été envoyé par Alice et non par
$\mathcal{A}$ se faisant passer pour elle~? La réponse est la signature
numérique asymétrique~: seul le détenteur de la clé privée d'Alice peut
produire sa signature. C'est le fondement mathématique du protocole ZK-NR
\cite{minka2025zknr} formalisé dans la Partie~VI.

\subsection{Du modèle Dolev-Yao au Trilemme CRO}

Dolev-Yao modélise ce que l'adversaire peut faire techniquement.
Le Trilemme CRO généralise en ajoutant la dimension légale~:

\begin{equation}
\underbrace{\mathcal{C}}_{\text{protéger contre }\mathcal{A}} +
\underbrace{\mathcal{R}}_{\text{résister à }\mathcal{A}} +
\underbrace{\mathcal{O}}_{\text{prouver face au juge}}
\;\leq\; 1 + \epsilon
\label{eq:CRO-dolev}
\end{equation}

La relation~\eqref{eq:CRO-dolev} formalise l'impossibilité~: maximiser la
protection contre $\mathcal{A}$ ($\mathcal{C}$ et $\mathcal{R}$ élevés) rend
la preuve difficile à vérifier par des tiers ($\mathcal{O}$ contraint).
Les protocoles ZK résolvent partiellement cette tension.

% ============================================================================
\section{Complexité Computationnelle et Sécurité Forensique}
\label{sec:complexite}
% ============================================================================

\subsection{Problèmes difficiles et sécurité cryptographique}

Toute la cryptographie moderne repose sur des problèmes
\textbf{calculatoirement difficiles}. La notation $\mathcal{O}$-grand
formalise la complexité~:

\begin{equation}
f(n) = \mathcal{O}(g(n)) \iff
\exists\, c > 0,\; n_0 \geq 0 \;\text{ tels que }\;
\forall n \geq n_0,\; f(n) \leq c\,g(n)
\label{eq:big-O}
\end{equation}

Les deux problèmes difficiles fondamentaux~:
\begin{enumerate}
    \item \textbf{Factorisation d'entiers}~: étant donné $N = p \cdot q$,
          retrouver $p$ et $q$. Complexité classique~:
          $\mathcal{O}(e^{n^{1/3}})$. Fondement de RSA.
    \item \textbf{Logarithme discret}~: étant donné $g^x \bmod p$,
          retrouver $x$. Fondement de Diffie-Hellman et des courbes
          elliptiques.
\end{enumerate}

\subsection{La menace quantique~: algorithme de Shor}

L'algorithme de \termecle{Shor}\index{Shor, algorithme de} (1994) résout
ces deux problèmes en temps polynomial sur un ordinateur quantique~:

\begin{equation}
T_{\text{Shor}}(n) = \mathcal{O}(n^3)
\quad \text{vs} \quad
T_{\text{classique}}(n) = \mathcal{O}(e^{n^{1/3}})
\label{eq:shor}
\end{equation}

L'implication pour la forensique est directe~: les preuves collectées en 2026
et chiffrées avec RSA ou ECDSA seront potentiellement déchiffrables dans
10 à 15 ans. C'est la stratégie
\termecle{Harvest Now, Decrypt Later}\index{Harvest Now Decrypt Later}~:
un adversaire étatique collecte aujourd'hui des données chiffrées pour les
déchiffrer demain.

\begin{boxalert}
\textbf{Migration post-quantique --- obligation éthique pour les données
longue durée.}

Le NIST a finalisé en 2024 quatre algorithmes post-quantiques~:
CRYSTALS-Kyber (chiffrement), CRYSTALS-Dilithium, FALCON et SPHINCS+
(signatures). Pour les investigations impliquant des données à longue
durée de vie (terrorisme, protection de témoins, secrets industriels),
le passage à ces algorithmes est une obligation éthique.
\end{boxalert}

\subsection{Fonctions à sens unique et hachage forensique}

Le hachage cryptographique repose sur une
\termecle{fonction à sens unique}\index{fonction à sens unique}~:

\begin{equation}
h : \{0,1\}^* \to \{0,1\}^n
\label{eq:hash-def}
\end{equation}

facile à calculer ($\mathcal{O}(|m|)$), impossible à inverser sans clé.
Trois propriétés de sécurité~:
\begin{enumerate}
    \item \textbf{Résistance à la préimage}~: étant donné $y$, infaisable
          de trouver $x$ tel que $h(x) = y$.
    \item \textbf{Résistance à la seconde préimage}~: étant donné $x$,
          infaisable de trouver $x' \neq x$ tel que $h(x') = h(x)$.
    \item \textbf{Résistance aux collisions}~: infaisable de trouver
          $(x, x')$ tels que $h(x) = h(x')$.
\end{enumerate}

Ces propriétés fondent le \textbf{principe de preuve d'intégrité}~: si le
hash de l'image forensique correspond au hash du disque original, aucune
modification n'a eu lieu. L'investigateur ne \emph{déclare} pas l'intégrité~;
le hash la \emph{prouve} mathématiquement.

% ============================================================================
\section{Modèles de Processus Forensique}
\label{sec:modeles-processus}
% ============================================================================

\subsection{Panorama et positionnement}

\begin{tableaucours}{p{2.4cm}p{2.0cm}p{5.5cm}p{1.8cm}}{%
    \tableauHeader{Modèle} &
    \tableauHeader{Origine} &
    \tableauHeader{Contribution} &
    \tableauHeader{Niveau}
}
DFRWS 2001 & Recherche & Premier modèle systématique~: 6~phases séquentielles & Conceptuel \\
Casey 2004 & Académique & Modèle intégré~: 5~phases, contexte judiciaire fort & Académique \\
\norme{NIST SP 800-86} & Gouvernement &
    Intégration IR + forensique~; distinction Exam./Analysis & Opérationnel \\
\norme{ISO/IEC 27037} & Normalisation &
    Standard international~: rôles DEFR et DES définis & Normatif \\
\end{tableaucours}
\vspace{-0.8em}
\begin{center}{\small\itshape Tableau~4.2 --- Panorama des modèles de processus}
\end{center}

\subsection{Le modèle DFRWS (2001)}

Le \termecle{DFRWS}\index{DFRWS} \cite{dfrws2001} a publié le premier modèle
formel de processus en six phases~:
\begin{enumerate}
    \item \textbf{Identification} --- reconnaître l'incident et définir son périmètre.
    \item \textbf{Préservation} --- isoler et protéger les preuves.
    \item \textbf{Collection} --- acquérir selon des méthodes documentées.
    \item \textbf{Examination} --- extraire les informations des preuves collectées.
    \item \textbf{Analysis} --- corréler pour reconstituer les événements.
    \item \textbf{Presentation} --- communiquer de façon défendable.
\end{enumerate}

La limite du DFRWS est son caractère \textbf{linéaire}~: un investigateur
alterne en réalité entre Collection et Examination au fil des découvertes.

\subsection{La distinction Examination / Analysis (NIST SP~800-86)}

La contribution centrale de \norme{NIST SP~800-86} \cite{nist80086} est la
distinction~:
\begin{itemize}
    \item \textbf{Examination}~: appliquer des méthodes systématiques pour
          identifier et extraire. L'examinateur documente ce qu'il \emph{trouve},
          pas ce que cela \emph{signifie}.
    \item \textbf{Analysis}~: interpréter les données extraites, tester des
          hypothèses, corréler les sources. L'analyste tire des conclusions.
\end{itemize}

Cette distinction est fondamentale~: un avocat peut attaquer une interprétation~;
il ne peut pas attaquer un hash qui correspond.

\subsection{ISO/IEC 27037~: standard international et rôles}

\norme{ISO/IEC~27037:2012} définit deux rôles distincts~:
\begin{itemize}
    \item \textbf{DEFR} (\textit{Digital Evidence First Responder}) ---
          premier intervenant sur la scène~; responsable de la préservation
          initiale. Son erreur est irréparable.
    \item \textbf{DES} (\textit{Digital Evidence Specialist}) --- analyste
          en laboratoire~; responsable de l'examination et de l'analysis.
\end{itemize}

\begin{boxjuris}
\textbf{Compatibilité ISO/IEC~27037 et droit camerounais~:}\\[0.3ex]
La \loicam{2010/012}, articles~52--59, impose~: présence des parties,
copie réalisée en leur présence, documentation. Ces exigences sont
\emph{substantiellement alignées} sur ISO/IEC~27037. Un investigateur
respectant le standard international respecte de facto les articles~52--59.
\end{boxjuris}

% ============================================================================
\section{Théorie de la Preuve Numérique}
\label{sec:theorie-preuve}
% ============================================================================

\subsection{Conditions formelles d'admissibilité}

La validité d'une preuve numérique $P$ peut être formalisée~:

\begin{equation}
\mathrm{Admissible}(P) \iff
\underbrace{\mathrm{Légale}(P)}_{\text{obtention licite}}
\;\wedge\;
\underbrace{\mathrm{Intègre}(P)}_{\text{hash vérifié}}
\;\wedge\;
\underbrace{\mathrm{Authentique}(P)}_{\text{custody continue}}
\;\wedge\;
\underbrace{\mathrm{Pertinente}(P)}_{\text{lien avec les faits}}
\label{eq:admissibilite}
\end{equation}

La conjonction ($\wedge$) est stricte~: l'absence d'une seule condition
invalide la preuve entière.

\subsection{La distinction constatation / interprétation}

\begin{tableaucours}{p{5.5cm}p{5.5cm}}{%
    \tableauHeader{Constatation (factuelle)} &
    \tableauHeader{Interprétation (analytique)}
}
« Le fichier \texttt{faux\_bon.docx} a été créé le 03/02/2026 à 14h22. » &
« BELLO a créé ce fichier pour préparer un détournement. » \\
« \texttt{svchost.exe} (PID~4892) a \texttt{cmd.exe} comme parent. » &
« Ce processus est anormal et indique une compromission. » \\
« Le hash MD5 de l'image correspond au hash de l'original. » &
« L'image est fidèle au disque et n'a pas été altérée. » \\
« L'adresse IP source est 185.220.XX.XX. » &
« L'attaque provenait de Roumanie. » \\
\end{tableaucours}
\vspace{-0.8em}
\begin{center}{\small\itshape Tableau~4.3 --- Distinction constatation/interprétation~:
exemples tirés des affaires TANGUE et MVOM}\end{center}

\subsection{La méthode des hypothèses concurrentes (ACH)}

Un investigateur rigoureux ne cherche pas à \emph{confirmer} sa thèse~:
il cherche à la \emph{falsifier}. Pour chaque conclusion~:

\begin{equation}
H_0 \;\text{(hypothèse principale)} \quad \text{vs} \quad
H_1, H_2, \ldots, H_k \;\text{(hypothèses concurrentes)}
\label{eq:hypotheses}
\end{equation}

Exemple dans MVOM~: $H_0$ = « Attaque par le groupe LockBit via RDP »..
Hypothèse $H_1$~: « La connexion RDP était légitime (télétravail autorisé)
et le ransomware a été introduit par un autre vecteur ». L'investigateur
doit \textbf{explicitement réfuter} $H_1$ --- en vérifiant la liste des
accès RDP autorisés, les contrats, les bulletins de salaire --- avant
de conclure à $H_0$.

\begin{boiteRemarque}{La méthode ACH (\textit{Analysis of Competing Hypotheses})}
Développée par Heuer pour la CIA \cite{heuer1999}, la méthode ACH formalise
cette approche~: pour chaque hypothèse, construire une matrice hypothèses/preuves
et calculer un score de cohérence. Une hypothèse qui explique toutes les preuves
et ne peut être réfutée par aucune est la conclusion la plus solide. Ce principe
est au cœur du Processus Investigatif Universel du
Chapitre~\ref{chap:PIU}.
\end{boiteRemarque}

% ============================================================================
\section{Synthèse~: Fondements Théoriques et Trilemme CRO}
\label{sec:fondements-cro}
% ============================================================================

\begin{tableaucours}{p{3.5cm}p{2.8cm}p{4.8cm}}{%
    \tableauHeader{Fondement} &
    \tableauHeader{Dim. CRO principale} &
    \tableauHeader{Contribution}
}
Locard numérique & $\mathcal{R}$ &
    Garantit l'existence de traces~; justifie l'exhaustivité de la collecte \\
Entropie de Shannon & $\mathcal{R}$ &
    Détecte les anomalies~; identifie chiffrement et stéganographie \\
Distance Hamming/Levenshtein & $\mathcal{R}$ &
    Mesure la similarité~; identifie variants et falsifications \\
Théorie des graphes & $\mathcal{O}$ &
    Visualise les relations~; cartographie le réseau criminel \\
Modèle Dolev-Yao & $\mathcal{C}$ &
    Modélise la menace adversariale~; justifie le chiffrement des preuves \\
Complexité / hash & $\mathcal{R} + \mathcal{O}$ &
    Fonde mathématiquement l'intégrité et l'authenticité \\
Théorie de la preuve & $\mathcal{O}$ &
    Formalise l'admissibilité et la séparation constat/interprétation \\
\end{tableaucours}
\vspace{-0.8em}
\begin{center}{\small\itshape Tableau~4.4 --- Fondements et dimensions CRO}
\end{center}

\begin{boxCRO}
\textbf{La tension fondamentale~:} maximiser $\mathcal{O}$ (transparence,
vérifiabilité) s'oppose à maximiser $\mathcal{C}$ (opacité, chiffrement).
La résolution partielle est apportée par les preuves à divulgation nulle~:
\textit{je prouve que je sais} sans révéler \textit{ce que je sais}.
Ce pont entre $\mathcal{C}$ et $\mathcal{O}$ est le sujet de la Partie~VI.
\end{boxCRO}

\separateur

\begin{boxresume}
\textbf{Ce qu'il faut retenir de ce chapitre~:}

\begin{itemize}
    \item \textbf{Locard numérique}~: quatre niveaux de traces (primaires,
          secondaires, tertiaires, résiduelles). L'anti-forensique laisse
          elle-même des traces secondaires.
    \item \textbf{Théorie de l'information}~: entropie ($H > 7.5$~: chiffré),
          fuzzy hashes (ssdeep) pour les variantes de malware, dnstwist pour
          le typosquatting.
    \item \textbf{Graphes}~: $G = (V,E,w)$ modélise les réseaux criminels~;
          la centralité d'intermédiarité identifie les ponts critiques~;
          les graphes temporels reconstituent les chemins d'attaque.
    \item \textbf{Dolev-Yao}~: adversaire omniscient sur le réseau, impuissant
          contre la cryptographie correcte. Justifie le chiffrement des preuves
          et la signature numérique.
    \item \textbf{Complexité}~: l'algorithme de Shor réduit RSA/ECDSA à
          $\mathcal{O}(n^3)$~; migration vers CRYSTALS-Dilithium/FALCON
          obligatoire pour les données longue durée.
    \item \textbf{Modèles de processus}~: DFRWS (6~phases), NIST SP~800-86
          (Exam./Analysis), ISO~27037 (DEFR/DES). La distinction
          constatation/interprétation est la règle d'or du rapport.
    \item \textbf{Admissibilité formelle}~:
          $\mathrm{Admissible}(P) \iff \mathrm{Légale} \wedge \mathrm{Intègre}
          \wedge \mathrm{Authentique} \wedge \mathrm{Pertinente}$~;
          la conjonction est stricte.
\end{itemize}
\end{boxresume}

\bigskip

\begin{boxexo}[title={\ding{45}\quad Exercice~4.1 --- Entropie et interprétation \niveaubase}]
Trois fichiers extraits d'une image forensique~:
\begin{enumerate}[(a)]
    \item \texttt{rapport\_finance.docx}~: entropie calculée = 4.7~bits/octet.
    \item \texttt{backup\_2026.zip}~: entropie calculée = 7.9~bits/octet.
    \item \texttt{systeme.dll}~: entropie calculée = 6.8~bits/octet.
\end{enumerate}
Pour chaque fichier~: interprétez l'entropie, indiquez si une investigation
approfondie est justifiée, et précisez ce que vous chercheriez.
\end{boxexo}

\begin{boxexo}[title={\ding{45}\quad Exercice~4.2 --- Graphe WOURI \niveaumoyen}]
À partir des éléments de l'Opération WOURI (quatre suspects, deux victimes,
flux de 47M et 32M~FCFA), construisez le graphe $G = (V, E, w)$~:
\begin{enumerate}[(a)]
    \item Définissez précisément $V$, $E$ et $w$.
    \item Calculez la centralité de degré de chaque suspect.
    \item Identifiez le suspect à la plus forte centralité d'intermédiarité
          et justifiez pourquoi sa neutralisation fragmente le réseau.
    \item Proposez quelle preuve supplémentaire permettrait d'ajouter une
          arête directe entre ATANGANA et HASSAN dans le graphe.
\end{enumerate}
\end{boxexo}

\begin{boxexo}[title={\ding{45}\quad Exercice~4.3 --- Admissibilité et ACH \niveauexpert}]
Dans l'affaire MVOM, l'investigateur a écrit dans son rapport~:

\medskip
\textit{« Le processus \texttt{svchost.exe} (PID~4892) était clairement
malveillant et a servi à l'installation du ransomware LockBit par les
attaquants russes. »}
\medskip

\begin{enumerate}[(a)]
    \item Identifiez au moins quatre problèmes dans cette phrase (distinctions
          constatation/interprétation, degré de certitude, attribution).
    \item Réécrivez-la en séparant strictement les constatations vérifiables
          des interprétations, avec un niveau de confiance explicite.
    \item Formulez la méthode ACH~: définissez $H_0$ et deux hypothèses
          concurrentes~; proposez quelle preuve discriminerait entre $H_0$
          et $H_1$.
\end{enumerate}
\end{boxexo}

\bigskip

\begin{boxinfo}
\textbf{Transition.} Le Chapitre~\ref{chap:etat-art} dresse l'état de l'art
de la discipline en 2025~: IA forensique, deepfake detection, forensique de
conteneurs, roadmap des certifications. Il complète la Partie~I avant le
plongeon opérationnel de la Partie~II, qui construit le Processus Investigatif
Universel sur les fondements posés ici.
\end{boxinfo}